% GENERATED FILE. Edit canonical lessons, then run npm run generate:books.
% canonical-source-hash: fnv1a64:7db63b3793dbd094
% canonical-lessons: KA-C189-gumpu, KA-C189-jatre, KA-C189-puje, KA-C189-devaru, KA-C189-upavasa

\chapter{Group, Village fair, Worship, God, Fasting}
\label{ch:ka-group-village-fair-worship-god-fasting}
\hlchaptermodality{189}

\begin{chapteropening}
\textbf{By the end of this chapter:} \emph{I can say a group, a village fair, worship, god and fasting.}

Complete the last lesson of chapter 189: I can say a group, a village fair, worship, god and fasting.
\end{chapteropening}

\section[\texorpdfstring{\kn{ಗುಂಪು} (gumpu)}{ಗುಂಪು (gumpu)}]{\kn{ಗುಂಪು} (gumpu) — a group, a crowd}

\label{lesson:KA-C189-gumpu}



\begin{quote}
Before the new one: say the Kannada for a thorn, then the Kannada for holy basil.
\end{quote}

\subsection*{You'll want to know: \kn{ಗುಂಪು}}
\textbf{\kn{ಗುಂಪು}} — \emph{gumpu} — \textquotedblleft{}a group, a crowd\textquotedblright{}.

\begin{grammarlens}[title={What you've built}]
One more word for this chapter.
\end{grammarlens}

\subsection*{Your turn}
\begin{itemize}
\raggedright
  \item \emph{Say it:} \emph{gumpu}
  \item \emph{Say it:} \emph{gumpu}, once more
  \item \emph{Say it:} say \emph{gumpu}
  \item \emph{Recall:} say the Kannada for a thorn, then the Kannada for holy basil, then say \emph{gumpu} again
\end{itemize}

\begin{tcolorbox}[breakable,colback=teal!4,colframe=teal!35!black,title={Before you move on}]
What does \kn{ಗುಂಪು} mean? (\textquotedblleft{}A group, a crowd\textquotedblright{}.) Say it once more.
\end{tcolorbox}
\section[\texorpdfstring{\kn{ಜಾತ್ರೆ} (jātre)}{ಜಾತ್ರೆ (jātre)}]{\kn{ಜಾತ್ರೆ} (jātre) — a village fair, a temple fair}

\label{lesson:KA-C189-jatre}



\begin{quote}
Before the new one: say the Kannada for holy basil, then the Kannada for a group.
\end{quote}

\subsection*{You'll want to know: \kn{ಜಾತ್ರೆ}}
\textbf{\kn{ಜಾತ್ರೆ}} — \emph{jātre} — \textquotedblleft{}a village fair, a temple fair\textquotedblright{}.

\begin{grammarlens}[title={What you've built}]
One more word for this chapter.
\end{grammarlens}

\subsection*{Your turn}
\begin{itemize}
\raggedright
  \item \emph{Say it:} \emph{jātre}
  \item \emph{Say it:} \emph{jātre}, once more
  \item \emph{Say it:} say \emph{jātre}
  \item \emph{Recall:} say the Kannada for holy basil, then the Kannada for a group, then say \emph{jātre} again
\end{itemize}

\begin{tcolorbox}[breakable,colback=teal!4,colframe=teal!35!black,title={Before you move on}]
What does \kn{ಜಾತ್ರೆ} mean? (\textquotedblleft{}A village fair, a temple fair\textquotedblright{}.) Say it once more.
\end{tcolorbox}
\section[\texorpdfstring{\kn{ಪೂಜೆ} (pūje)}{ಪೂಜೆ (pūje)}]{\kn{ಪೂಜೆ} (pūje) — worship, a puja}

\label{lesson:KA-C189-puje}



\begin{quote}
Before the new one: say the Kannada for a group, then the Kannada for a village fair.
\end{quote}

\subsection*{You'll want to know: \kn{ಪೂಜೆ}}
\textbf{\kn{ಪೂಜೆ}} — \emph{pūje} — \textquotedblleft{}worship, a puja\textquotedblright{}.

\begin{grammarlens}[title={What you've built}]
One more word for this chapter.
\end{grammarlens}

\subsection*{Your turn}
\begin{itemize}
\raggedright
  \item \emph{Say it:} \emph{pūje}
  \item \emph{Say it:} \emph{pūje}, once more
  \item \emph{Say it:} say \emph{pūje}
  \item \emph{Recall:} say the Kannada for a group, then the Kannada for a village fair, then say \emph{pūje} again
\end{itemize}

\begin{tcolorbox}[breakable,colback=teal!4,colframe=teal!35!black,title={Before you move on}]
What does \kn{ಪೂಜೆ} mean? (\textquotedblleft{}Worship, a puja\textquotedblright{}.) Say it once more.
\end{tcolorbox}
\section[\texorpdfstring{\kn{ದೇವರು} (dēvaru)}{ದೇವರು (dēvaru)}]{\kn{ದೇವರು} (dēvaru) — god}

\label{lesson:KA-C189-devaru}



\begin{quote}
Before the new one: say the Kannada for a village fair, then the Kannada for worship.
\end{quote}

\subsection*{You'll want to know: \kn{ದೇವರು}}
\textbf{\kn{ದೇವರು}} — \emph{dēvaru} — \textquotedblleft{}god\textquotedblright{}.

\begin{grammarlens}[title={What you've built}]
One more word for this chapter.
\end{grammarlens}

\subsection*{Your turn}
\begin{itemize}
\raggedright
  \item \emph{Say it:} \emph{dēvaru}
  \item \emph{Say it:} \emph{dēvaru}, once more
  \item \emph{Say it:} say \emph{dēvaru}
  \item \emph{Recall:} say the Kannada for a village fair, then the Kannada for worship, then say \emph{dēvaru} again
\end{itemize}

\begin{tcolorbox}[breakable,colback=teal!4,colframe=teal!35!black,title={Before you move on}]
What does \kn{ದೇವರು} mean? (\textquotedblleft{}God\textquotedblright{}.) Say it once more.
\end{tcolorbox}
\section[\texorpdfstring{\kn{ಉಪವಾಸ} (upavāsa)}{ಉಪವಾಸ (upavāsa)}]{\kn{ಉಪವಾಸ} (upavāsa) — fasting, a fast}

\label{lesson:KA-C189-upavasa}



\begin{quote}
Before the new one: say the Kannada for worship, then the Kannada for god.
\end{quote}

\subsection*{You'll want to know: \kn{ಉಪವಾಸ}}
\textbf{\kn{ಉಪವಾಸ}} — \emph{upavāsa} — \textquotedblleft{}fasting, a fast\textquotedblright{}.

\begin{grammarlens}[title={What you've built}]
One more word for this chapter.
\end{grammarlens}

\subsection*{Your turn}
\begin{itemize}
\raggedright
  \item \emph{Say it:} \emph{upavāsa}
  \item \emph{Say it:} \emph{upavāsa}, once more
  \item \emph{Say it:} say \emph{upavāsa}
  \item \emph{Recall:} say the Kannada for worship, then the Kannada for god, then say \emph{upavāsa} again
\end{itemize}

\begin{tcolorbox}[breakable,colback=teal!4,colframe=teal!35!black,title={Before you move on}]
What does \kn{ಉಪವಾಸ} mean? (\textquotedblleft{}Fasting, a fast\textquotedblright{}.) Say it once more.
\end{tcolorbox}
